% Exercises \proofgraphlabel, checked against reference.dot:
%   a format using the number alone after a letter (L#1);
%   a format with a tie, written out as a space (Prop.~#1);
%   a format using the printed name as well (#2 (#1));
%   an environment left alone, which keeps the default label;
%   an unnumbered environment drawn through \proofgraphtrack, whose #1 is empty.
% And \proofgraphnodelabel, a label of its own for one result:
%   over the default label, with an accented letter;
%   over a \proofgraphlabel format, with quotes, which the .dot has to
%   escape;
%   the same text for two results, which remain two nodes with their edges.
\documentclass{article}
\usepackage{amsthm}
\usepackage{proofgraph}
\newtheorem{theorem}{Theorem}
\newtheorem{lemma}{Lemma}
\newtheorem{proposition}{Proposition}
\newtheorem{corollary}{Corollary}
\newtheorem*{mainthm}{Main Theorem}
\proofgraphtrack{mainthm}
\proofgraphlabel{lemma}{L#1}
\proofgraphlabel{proposition}{Prop.~#1}
\proofgraphlabel{corollary}{#2 (#1)}
\proofgraphlabel{mainthm}{Main#1}
\begin{document}
\begin{lemma}\label{l}A lemma.\end{lemma}
\begin{proposition}\label{p}A proposition.\end{proposition}
\begin{proof}By \ref{l}.\end{proof}
\begin{theorem}\label{t}A theorem.\end{theorem}
\begin{proof}By \ref{p}.\end{proof}
\begin{corollary}\label{c}A corollary.\end{corollary}
\begin{proof}By \ref{t}.\end{proof}
\begin{mainthm}\label{m}The main theorem.\end{mainthm}
\begin{proof}By \ref{c}.\end{proof}
\begin{theorem}\label{r}\proofgraphnodelabel{Théorème de Rolle}
A named theorem.\end{theorem}
\begin{proof}By \ref{m}.\end{proof}
\begin{lemma}\label{q} \proofgraphnodelabel{The "key" lemma}
A lemma with a format, named all the same.\end{lemma}
\begin{proof}By \ref{r}.\end{proof}
\begin{lemma}\label{d1}\proofgraphnodelabel{Twin}A twin.\end{lemma}
\begin{proof}By \ref{q}.\end{proof}
\begin{lemma}\label{d2}\proofgraphnodelabel{Twin}Its twin.\end{lemma}
\begin{proof}By \ref{r}.\end{proof}
\end{document}
